\section{An interior injectivity theorem}
\label{sec:interior-injectivity}

We prove \Cref{thm:interior-injectivity} from its two endpoint
metrics, using the top-degree $L^2$ estimate in
\Cref{lem:top-degree-l2}. The proof does not use
\Cref{thm:adjoint-metric} or its extremal-section construction. The output is the natural map on ordinary coherent cohomology.
The proof therefore needs both a weighted estimate on the complement
of the crossing divisor and a comparison that retains the original
ordinary class. The fixed-cover construction below supplies this
comparison together with a uniform bound for primitives.

For a class in the kernel, we will construct weighted harmonic
representatives whose images have norms tending to zero. The endpoint
metrics make multiplication a contraction and supply the curvature
comparison for this estimate. Local solutions of the $\dbar$ equation
leave holomorphic differences on overlaps; controlling a splitting of
these differences on one fixed cover gives uniformly bounded global
primitives. The extended differences form a cocycle representing the
ordinary cohomology class of the original form. Finally, comparison in
one fixed Hilbert space lets the small images force the original
ordinary class to vanish.
Classical injectivity for roots of divisors, including an additional
reduced boundary, is proved in \cite[Theorem~5.1]{EsnaultViehweg1992}.
The local and uniform-primitive methods have precedents in
\cite[Sections~5.2--5.3]{Matsumura2018}; the precise comparison needed
for the present line bundles is proved here.

\subsection{Endpoint regularization and the comparison weights}

Fix a K\"ahler form $\omega$ on $V$, and put $n=\dim V$.  The
zero-dimensional assertion is trivial, so assume $n>0$.  We use the
following consequence of regularization and exponential integrability.

\begin{lemma}\label[lemma]{lem:zero-lelong-regularization}
Let $F$ be a rational line bundle on a smooth projective complex
variety $V$ equipped with a K\"ahler form $\omega$, and suppose that
$F$ has a semipositive singular metric $\psi$ having zero Lelong
numbers everywhere.  There are smooth weights $\psi_k$ on $F$ and
numbers $\epsilon_k\downarrow0$ such that
\[
 i\partial\bar\partial\psi_k\ge-\epsilon_k\omega.
\]
In fixed local frames these weights are uniformly bounded above on
relatively compact coordinate neighborhoods, and, for every fixed
$b>0$, their exponentials $e^{-b\psi_k}$ have uniformly bounded local
$L^1$ norms.
\end{lemma}

\begin{proof}
Choose a smooth reference metric on $F$. Apply the attenuation form of
Demailly's regularization theorem to the specified current
$T=i\partial\bar\partial\psi$, using the global difference between
$\psi$ and the reference weight as its potential
\cite[Main Theorem~1.1]{Demailly1992}. Fix a positive attenuation
level $c$. Adding the reference weight back to the approximating
potentials gives weights $\psi_k$ on the same rational bundle, with
$\psi_k\downarrow\psi$. They are smooth outside the Lelong level set
$E_c(T)=\{x:\nu(T,x)\ge c\}$, which is empty. Their curvature lower bounds have the form
$-\min(\lambda_k,c)u-\delta_k\omega$, where $u$ is one fixed
smooth nonnegative form, $\delta_k\to0$, and the continuous functions
$\lambda_k$ decrease to the Lelong-number function. That function
is zero here. Dini's theorem on the compact manifold gives
$\sup_V\lambda_k\to0$, and $u\le C\omega$ gives the stated
loss $\epsilon_k\to0$. Passing to a subsequence makes the losses
decrease. Work first on a common integral tensor power of $F$ and
divide all weights by that integer afterwards.

On each fixed coordinate neighborhood the decreasing approximants
satisfy $\psi\le\psi_k\le\psi_1$. The upper bound is smooth on a
slightly larger neighborhood, and hence is bounded on the chosen
compact subset. For every fixed $b>0$,
\[
                      e^{-b\psi_k}\le e^{-b\psi}.
\]
The right side is locally integrable by the zero-Lelong case of Skoda's
integrability criterion \cite{Skoda1972}; see
\cite[Property~1.4(8)]{DemaillyKollar2001} for the formulation for
arbitrary plurisubharmonic weights. This proves the uniform
fixed-exponent bounds directly.
\end{proof}

Apply the lemma simultaneously to $L-C_0$ and $L-C_2$, obtaining
weights $\psi_{0,k},\psi_{2,k}$ with a common sequence
$0<\epsilon_k\le1$ tending to zero.  Let
\[
 D=\Supp(C_0+C_2)_{\rm red},\qquad U=V\setminus D.
\]
For a rational divisor $C$, write $[C]$ for its divisor weight.
The endpoint weights $\psi_{i,k}$ live on the rational bundles
$L-C_i$. Adding $[C_i]$ gives weights on $L$. We first raise the
weight at the first endpoint above that at the second, so that the
natural inclusion of the eventual integral bundles will be a
contraction. The symbols $\Phi_{i,k}$ below are weights on $L$,
whereas $h_{i,k}$ are weights on $L_i=L-\lfloor C_i\rfloor$.
On $U$ set
\begin{align}
 \psi'_{0,k}
   &=\log\bigl(e^{\psi_{0,k}}
             +e^{\psi_{2,k}+[C_2-C_0]}\bigr),
       \label{eq:logsum-endpoint}\\
 \Phi_{0,k}&=\psi'_{0,k}+[C_0],\qquad
 \Phi_{2,k}=\psi_{2,k}+[C_2],\qquad
 \Phi_{1,k}=(1-\lambda)\Phi_{0,k}+\lambda\Phi_{2,k},
       \label{eq:interior-weights}\\
 h_{i,k}&=\Phi_{i,k}-[\floor{C_i}],\qquad i=0,1.
       \label{eq:integral-line-weights}
\end{align}
The two summands in \Cref{eq:logsum-endpoint} are weights on the same
rational line bundle $L-C_0$, so the expression is intrinsically
well-defined.  The logarithm of a sum of exponentials preserves the
common curvature lower bound: locally add a potential for
$\epsilon_k\omega$ to both summands and use the psh log-sum
inequality.  Moreover $\psi'_{0,k}\ge\psi_{0,k}$.  Effectivity of
$C_2-C_0$ gives a uniform local upper bound for its divisor weight, so
$\psi'_{0,k}$ retains both the upper bounds and all the fixed-exponent
integrability bounds of \Cref{lem:zero-lelong-regularization}.

The weights $h_{i,k}$ are smooth metrics on $L_i|_U$.  They have
uniform local upper bounds in frames extending across $D$, and
\begin{equation}\label{eq:weight-integral-bounds}
 \sup_k\int_{B\cap U}e^{-h_{i,k}}\,d\lambda_B<\infty
 \quad(i=0,1)
\end{equation}
on every relatively compact coordinate neighborhood $B$, where
$d\lambda_B$ is coordinate volume.  To verify the last assertion,
write
\[
 h_{0,k}=\psi'_{0,k}+[C_0-\floor{C_0}],\qquad
 h_{1,k}=(1-\lambda)\psi'_{0,k}+\lambda\psi_{2,k}
                      +[C_1-\floor{C_1}].
\]
Every coefficient of either fractional divisor lies in $[0,1)$.
Choose one H\"older exponent $a>1$ so close to one that its product
with each of these finitely many coefficients is still less than one.
The exponential of the fractional divisor weight is in $L^a$ by the
crossing-coordinate integral.  The conjugate H\"older exponent for
the remaining weights is allowed by the fixed-exponent bounds; a
second H\"older inequality handles the two summands in $h_{1,k}$.
This proves \Cref{eq:weight-integral-bounds} uniformly in $k$.

Let $E=\floor{C_1}-\floor{C_0}\ge0$, and let
$s:L_1\to L_0$ be multiplication by the canonical section of $E$.
Since $\Phi_{2,k}\le\Phi_{0,k}$, we have
$\Phi_{1,k}\le\Phi_{0,k}$ and
\begin{equation}\label{eq:multiplication-contraction}
 |s|^2e^{h_{1,k}-h_{0,k}}
 =e^{\Phi_{1,k}-\Phi_{0,k}}\le1.
\end{equation}
Thus multiplication is a contraction on all form degrees for any one
fixed background metric.  On $U$ the curvatures are
$\theta_{i,k}=i\partial\bar\partial\Phi_{i,k}$; divisor weights
contribute no curvature there.

\subsection{A complete metric with bounded local potentials}

The singular set must be removed to use smooth weighted harmonic
forms.  We choose the complete metric carefully so that local
solvability still has a uniform comparison with the given weights.
The bounded-potential construction follows the complete-metric method
of \cite[Lemma~3.1]{FujinoInjectivity2012}; we include the calculation
because both bounded potentials and completeness are needed in the
uniform primitive estimate.

\begin{lemma}\label{lem:bounded-potential-complete-metric}
There is a complete K\"ahler metric $\omega_c\ge\omega$ on $U$ such
that, on sufficiently small coordinate neighborhoods in $V$, it has
potentials bounded even on approach to $D$.
\end{lemma}

\begin{proof}
If $D$ is empty, take $\omega_c=\omega$.  Otherwise choose a smooth
metric on $\cO_V(D)$ and scale it so that
$t=\log|s_D|_{\rm sm}^2<-e^2$ on $U$.  Set $f_0(t)=1/\log(-t)$.
For $x=-t$,
\[
 f_0'(t)=\frac1{x(\log x)^2},\qquad
 f_0''(t)=\frac1{x^2}\left(\frac1{(\log x)^2}
                          +\frac2{(\log x)^3}\right).
\]
In particular $f_0'$ is bounded and
$f_0''(t)\ge(t\log(-t))^{-2}$.  Off $D$ the form
$i\partial\bar\partial t$ is the negative of a fixed smooth
curvature form, so it is bounded by a multiple of $\omega$.  For a
sufficiently large fixed $M$ the form
\[
 \omega_c=M\omega+i\partial\bar\partial f_0(t)
\]
is positive and dominates both $\omega$ and
$i\partial t\wedge\bar\partial t/(t\log(-t))^2$.

A path leaving every compact subset of $U$ approaches $D$, because
$V$ is compact; along such an approach $t\to-\infty$.  The displayed
radial term bounds its length below by a positive constant times the
total variation of $\log\log(-t)$, which diverges.  Hence the metric
is complete, including at crossings of components of $D$.  Finally,
if $\omega=i\partial\bar\partial\rho$ on a small coordinate
neighborhood, then $M\rho+f_0(t)$ is a local potential for $\omega_c$.
The smooth function $\rho$ is bounded on a smaller relatively compact
neighborhood, and $f_0(t)$ is bounded and tends to zero at $D$.
\end{proof}

Fix this metric for the rest of the section. When $D\ne\varnothing$,
the function $r=\log\log(-t)$ is a smooth proper exhaustion of $U$
with bounded gradient in $\omega_c$. If $\zeta$ is a smooth function
equal to one on $(-\infty,1]$ and zero on $[2,\infty)$, then
$\zeta(r/R)$, as $R\to\infty$, gives compactly supported smooth
cutoffs tending to one and with gradient $O(R^{-1})$. When $D$ is
empty, the constant cutoff one suffices.

The curvature comparison
following \Cref{eq:interior-weights} becomes
\begin{equation}\label{eq:curvature-interior-comparison}
 \theta_{1,k}+\epsilon_k\omega_c
 \ge(1-\lambda)(\theta_{0,k}+\epsilon_k\omega_c)\ge0,
\end{equation}
since $\theta_{2,k}+\epsilon_k\omega_c\ge0$ as well.

\subsection{A fixed-cover primitive estimate and the ordinary class map}

Fix a holomorphic line bundle $F$ on $V$. The local $L^2$ theorem
gives primitives on coordinate neighborhoods,
but a global primitive requires their holomorphic differences on
overlaps to split. We use one finite cover and its fixed spaces of
holomorphic sections to make that splitting estimate uniform in the
weights. The extended differences also define a class in the ordinary
group $H^1(V,K_V\otimes F)$. Thus the class map below has ordinary
coherent cohomology as its target, and the primitive estimate shows
that its kernel consists of actual weighted exact forms.
This local-solutions and holomorphic-cocycle comparison follows the
method of \cite[Section~3, Claim~1 and Lemma~3.2]{FujinoInjectivity2012}.
For a smooth metric $h$ on
$F|_U$, write $L^2_h(n,q)$ for the Hilbert space of $F$-valued
$(n,q)$ forms using $h$ and $\omega_c$.  The Dolbeault operator is its
maximal distributional realization, and
$Z^1_h=\ker(\dbar:L^2_h(n,1)\to L^2_h(n,2))$.

\begin{proposition}\label[proposition]{prop:uniform-primitives}
Let $F$ be a fixed holomorphic line bundle on $V$, and let $h_k$ be
smooth metric weights on $F|_U$.  Suppose that, for a fixed constant
$A_0\ge0$,
\[
 i\partial\bar\partial h_k\ge-A_0\omega_c,
\]
and that in local frames on $V$ the weights are uniformly bounded
above and the functions $e^{-h_k}$ have uniformly bounded local
integrals on punctured coordinate neighborhoods.  Then:
\begin{enumerate}[label=\textup{(\roman*)}]
\item There is a continuous linear map
\[
 \kappa_k:Z^1_{h_k}\longrightarrow H^1(V,K_V\otimes F)
\]
agreeing with the ordinary Dolbeault class on smooth closed forms on
$V$.  It annihilates the closure in $L^2_{h_k}(n,1)$ of the image of
$\dbar:L^2_{h_k}(n,0)\to L^2_{h_k}(n,1)$.
\item There is a constant $C$ independent of $k$ such that every
$f\in Z^1_{h_k}$ with $\kappa_k(f)=0$ has a global primitive on $U$
satisfying
\[
 \dbar v=f,\qquad \|v\|_{h_k}\le C\|f\|_{h_k}.
\]
No continuous linear choice of these global primitives is asserted.
\end{enumerate}
\end{proposition}

\begin{proof}
Choose a finite collection of coordinate and frame balls
\[
 V_j\Subset U_j\Subset\widehat U_j
\]
such that the $V_j$ still cover $V$ and $\omega_c$ has a bounded local
potential $\rho_j$ on $\widehat U_j\cap U$.  These sets and potentials
are independent of $k$.

\emph{Local solutions.}
On $U_j\cap U$, replace the line weight by
$h_k+(A_0+1)\rho_j$.  Its curvature dominates $\omega_c$, and the
bounded potential compares its top-form norms with those of $h_k$ by
fixed constants.  Equip $U_j\cap U$ with the complete metric
$\omega_c+\omega_{\rm ball}$, where $\omega_{\rm ball}$ is a complete
metric for the coordinate ball.  This sum is complete at both the
deleted divisor and the ball boundary.  An $(n,1)$ datum square
integrable for $\omega_c$ remains so for the larger background metric:
its density uses the inverse background matrix on its one remaining
covector.  By \Cref{lem:top-degree-l2}, the inverse-curvature density
for the modified weight is bounded by the original $\omega_c$ datum
density.  The top-form solution density is independent of the
auxiliary complete metric.  Consequently there is a local solution
operator $T_{j,k}$ with
\begin{equation}\label{eq:local-primitive-bound}
 \dbar T_{j,k}f=f,\qquad
 \|T_{j,k}f\|_{h_k,U_j\cap U}
       \le C_j\|f\|_{h_k,U_j\cap U},
\end{equation}
where $C_j$ is independent of $k$.  For each fixed $k$, choose the
solution of smallest norm in the modified metric.  Orthogonal
projection onto the complement of the holomorphic kernel makes this
a linear continuous choice.  The affine solution set is closed because
the distributional operator is closed.

\emph{Holomorphic cocycles.}
Given $f\in Z^1_{h_k}$, put $w_j=T_{j,k}(f|_{U_j\cap U})$.
On $(U_i\cap U_j)\cap U$, the difference $w_j-w_i$ is holomorphic.
Uniform upper bounds for $h_k$ turn
\Cref{eq:local-primitive-bound} into ordinary local $L^2$ bounds in
frames on $V$.  Holomorphic $L^2$ removal therefore extends the
differences uniquely across $D$.  One can see the removal locally by
Laurent expansion in the crossing coordinates and Fubini: a nonzero
negative power is not square integrable.  The extended differences
$c_{ij}=w_j-w_i$, sections of $K_V\otimes F$ on $U_i\cap U_j$, form a holomorphic
\v Cech cocycle.

For later use, this family of cocycles is bounded in every fixed
compact-convergence seminorm whenever $\|f\|_{h_k}$ is bounded,
uniformly in $k$.  Indeed, for any compact set in an overlap, choose a
slightly larger compact neighborhood still in that overlap.  The
ordinary $L^2$ bound there and the holomorphic mean-value inequality
bound the supremum on the smaller compact.  Each such compact has
positive separation from the ball boundaries.  Its constant may
depend on that compact, but not on $k$.  No bound on the supremum over
an entire open overlap is required.

\emph{The ordinary class.}
Let $\mathcal Z^1$ be the Fr\'echet space of holomorphic cocycles on
this finite cover, with compact-convergence topology.  The cocycle
equations make it a closed subspace of the finite product of overlap
section spaces.  Choose a smooth partition of unity $(\eta_i)$ with
$\Supp\eta_i\Subset U_i$.  For a cocycle $c$ put
$b_j=\sum_i\eta_i c_{ij}$ on $U_j$, extending each summand by zero
outside its overlap.  Then $b_j-b_i=c_{ij}$, so the forms $\dbar b_j$
agree and define a global smooth form $\beta_c$.  This defines the
continuous linear map
\[
 \kappa:\mathcal Z^1\longrightarrow H^1(V,K_V\otimes F),
 \qquad c\longmapsto[\beta_c].
\]
Continuity follows from the compact support of the partition and
Cauchy estimates for derivatives on slightly larger compact sets;
ordinary cohomology is finite-dimensional and Hausdorff on the compact
manifold.  Define $\kappa_k(f)=\kappa((c_{ij}))$.
Changing the local primitives
changes the cocycle by a holomorphic coboundary: the differences of
two primitives extend by the same ordinary $L^2$ argument.  The local
solution estimates and the compact-seminorm estimates just proved
make $\kappa_k$ continuous and linear.

If $f$ is smooth on $V$, choose smooth local primitives on the larger
balls.  On the smaller balls they have finite weighted norms, because
their coefficients are bounded and the local integrals of $e^{-h_k}$
are finite.  Comparison with these primitives shows that $\kappa_k$
agrees with the ordinary Dolbeault class.  If instead $f=\dbar v$ for
a global weighted $L^2$ primitive on $U$, then $w_j-v$ is holomorphic
on $U_j\cap U$ and ordinary locally $L^2$.  Its extension gives a
holomorphic splitting of $c_{ij}$, so $\kappa_k(f)=0$.  Continuity
therefore annihilates the closure of the actual exact forms, as in
(i).

\emph{A uniformly bounded splitting.}
Let
$\mathcal C^0=\prod_jH^0(U_j,K_V\otimes F)$, again with
compact-convergence topology, and let $\mathcal B^1$ be the image of
its coboundary map $\delta$ in $\mathcal Z^1$.  We need the equality
\begin{equation}\label{eq:closed-coboundary-kernel}
 \mathcal B^1=\ker\kappa.
\end{equation}
The forward inclusion is immediate.  For the reverse inclusion,
construct smooth local cochains from a partition of unity for a
cocycle in $\ker\kappa$.  Their common $\dbar$ is a global smooth
form with zero ordinary Dolbeault class.  Subtract a global smooth
primitive from every cochain.  The resulting cochains are
holomorphic and split the original cocycle.  This proves
\Cref{eq:closed-coboundary-kernel} without any assumption that the
coordinate cover is a Leray cover.

It follows that $\mathcal B^1$ is a closed Fr\'echet subspace.  The
continuous surjection
$\delta:\mathcal C^0\to\mathcal B^1$ is open.  Prescribe only the
single continuous seminorm
\[
 q((a_j))=\max_j\sup_{\overline{V_j}}|a_j|,
\]
where the norms are those of fixed smooth frames or a fixed smooth
bundle metric.  Openness says that the image of $\{q<1\}$ contains a
zero neighborhood of $\mathcal B^1$ described by finitely many compact
seminorms.  The cocycles arising from all $k$ and all data of norm at
most one are bounded in those seminorms.  A single rescaling therefore
gives, for every such cocycle in $\mathcal B^1$, a splitting
$(a_j)$ with $q((a_j))\le C_{\rm split}$ for a fixed
$C_{\rm split}$.  This uses no
assertion that a bounded subset has a lift bounded in every
Fr\'echet seminorm.

When $\kappa_k(f)=0$, choose this splitting and set
$v=w_j-a_j$ on $U_j\cap U$.  The expressions agree on the full
overlaps, and hence define a global primitive on $U$.  Estimate its
norm on the smaller $V_j$, which cover $V$.  The local primitives
are controlled by \Cref{eq:local-primitive-bound}; the corrections
are controlled by their supremum on $\overline{V_j}$ and the uniform
integral of $e^{-h_k}$ there.  Summing over the finite cover proves a
uniform bound for unit-norm data.  Scaling proves (ii).
\end{proof}

Both sequences $h_{0,k}$ and $h_{1,k}$ satisfy this proposition, by
\Cref{eq:weight-integral-bounds,eq:curvature-interior-comparison}.
For each weight, parts (i) and (ii) give
\[
 \ker\kappa_k=\im\dbar=\overline{\im\dbar}
 \quad\text{inside }Z^1_{h_k}.
\]
Indeed, part (i) kills the closure of the exact forms, while part (ii)
provides a primitive for every form in the kernel. We will keep track
of the closures explicitly when comparing different weighted spaces.
The proposition also applies to a fixed smooth line metric from $V$: its curvature
is bounded below by $-A_0\omega_c$ for some fixed $A_0$, and its
local upper and integral bounds are automatic.  This last observation
will allow us to compare all approximation parameters in one Hilbert
space.

\subsection{Harmonic representatives and the Bochner comparison}

Take a class in the kernel of the map in
\Cref{thm:interior-injectivity}, and represent it by a smooth closed
$L_1$-valued $(n,1)$ form $\gamma$ on $V$.  All norms in this
subsection use $\omega_c$; a subscript $i,k$ indicates the line
weight $h_{i,k}$, and integrals of pointwise inner products use
$dV_{\omega_c}$.  The norms $\|\gamma\|_{1,k}$ are bounded
uniformly.  Indeed, enlarging $\omega$ to $\omega_c$ decreases the
top-form-valued $(0,1)$ density, and the coefficients of $\gamma$ are
bounded on a finite cover, where \Cref{eq:weight-integral-bounds}
applies.

In the maximal $L^2_{1,k}$ Dolbeault complex, the closure of the image
of degree-zero forms lies in the closed kernel of the degree-one
operator.  Subtract its orthogonal projection from $\gamma$, obtaining
$u_k$.  Thus
\begin{equation}\label{eq:reduced-harmonic-representative}
 \gamma-u_k\in\overline{\im\dbar},\qquad
 \dbar u_k=0,\qquad \dbar^*_{1,k}u_k=0,\qquad
 \|u_k\|_{1,k}\le\|\gamma\|_{1,k}.
\end{equation}
The adjoint assertion follows from orthogonality to the actual range;
it does not require that range to be closed.  Local elliptic
regularity for the smooth weighted Dolbeault Laplacian on $U$ makes
$u_k$ smooth there.

Write $\Lambda=\Lambda_{\omega_c}$ and let $D'_{i,k}$ be the
$(1,0)$ Chern derivative.  In bidegree $(n,1)$ put
\[
 A_{i,k}=[\theta_{i,k}+\epsilon_k\omega_c,\Lambda].
\]
These are pointwise nonnegative operators, and
$A_{1,k}\ge(1-\lambda)A_{0,k}$.  The smooth compact-support
Bochner--Kodaira identity is
\begin{equation}\label{eq:bochner-top-degree}
 \|\dbar v\|_{i,k}^2+\|\dbar^*_{i,k}v\|_{i,k}^2
 =\|D_{i,k}^{\prime *}v\|_{i,k}^2
       +\int_U\langle[\theta_{i,k},\Lambda]v,v\rangle_{i,k}.
\end{equation}
Here $D'_{i,k}v=0$ in top holomorphic degree, and
$[\omega_c,\Lambda]$ is the identity in this bidegree.  The same
identity on complete manifolds, with the indicated domain
interpretation and lower curvature bound, is discussed in
\cite[Proposition~2.4]{Matsumura2018}.

To justify its use without already knowing the extra terms are
integrable, use the compactly supported smooth cutoffs constructed above, whose
gradients tend uniformly to zero in $\omega_c$.  Apply
\Cref{eq:bochner-top-degree} to the cutoff multiples of $u_k$ and add
$\epsilon_k$ times their squared norms.  The left-hand derivative
terms tend to zero by \Cref{eq:reduced-harmonic-representative} and
the gradient bound.  The two terms retained on the right are
nonnegative, so lower semicontinuity gives
\begin{equation}\label{eq:source-bochner-energy}
 \|D_{1,k}^{\prime *}u_k\|_{1,k}^2
   +\int_U\langle A_{1,k}u_k,u_k\rangle_{1,k}
 \le\epsilon_k\|u_k\|_{1,k}^2.
\end{equation}
The commutators with the cutoffs contain only their derivatives; they
do not introduce constants from derivatives of the varying line
metrics.

The K\"ahler identity
$D^{\prime *}=i[\Lambda,\dbar]$ shows that this formal operator
commutes with holomorphic multiplication by $s$.  Combining
\Cref{eq:multiplication-contraction,eq:curvature-interior-comparison}
with \Cref{eq:source-bochner-energy} therefore gives
\begin{equation}\label{eq:target-bochner-energy}
 \|D_{0,k}^{\prime *}(su_k)\|_{0,k}^2
 +\int_U\langle A_{0,k}su_k,su_k\rangle_{0,k}
 \le C_\lambda\epsilon_k\|u_k\|_{1,k}^2.
\end{equation}
For the curvature term one first uses
$A_{1,k}\ge(1-\lambda)A_{0,k}$ and then the line-norm contraction.
One may take $C_\lambda=(1-\lambda)^{-1}$: this factor bounds both
the derivative term and the curvature term in the source energy.
The strict interior hypothesis keeps this comparison factor finite.

We also need the target $\dbar$ adjoint.  The form $su_k$ is smooth,
closed, and square integrable.  Apply the compact-support identity to
cutoff multiples of this form.  On its right-hand side use
\Cref{eq:target-bochner-energy}, replace the curvature commutator by
the larger nonnegative $A_{0,k}$, and let the gradient errors tend to
zero.  It follows that the formal adjoint is square integrable and
\begin{equation}\label{eq:small-target-adjoint}
 \|\dbar^*_{0,k}(su_k)\|_{0,k}
       \le C\sqrt{\epsilon_k}.
\end{equation}
This also verifies the Hilbert-adjoint domain condition.  In detail,
complete-metric cutoffs and local smooth approximation give density in
the maximal graph norm, so integration by parts against compactly
supported forms extends to that domain.  The formal adjoint, now
known to be $L^2$, is its Hilbert adjoint.  Thus
\Cref{eq:small-target-adjoint} has not assumed the existence of a
global primitive or an adjoint-domain statement in advance.

We have shown that the image of the harmonic representative has a
small adjoint norm.  The fixed-cover proposition supplies a uniformly
bounded primitive for this image; pairing these two statements will
make its entire target norm tend to zero.

\subsection{From a small image to vanishing of the ordinary class}

Multiplication by $s$ is bounded in degrees zero and one and commutes
with the distributional Dolbeault operator.  Hence
\Cref{eq:reduced-harmonic-representative} implies that
$s\gamma-su_k$ is in the target closure of the actual exact forms.
Apply the target class map in \Cref{prop:uniform-primitives}.  It
annihilates that closure and agrees with ordinary cohomology on
$s\gamma$.  The latter class is zero by the choice of $\gamma$.
Consequently $\kappa_{0,k}(su_k)=0$.

The norms of $su_k$ are uniformly bounded by
\Cref{eq:multiplication-contraction,eq:reduced-harmonic-representative}.
Part (ii) of \Cref{prop:uniform-primitives} gives global forms $v_k$
on $U$ with
\[
 \dbar v_k=su_k,\qquad \sup_k\|v_k\|_{0,k}<\infty.
\]
The adjoint-domain conclusion above now justifies
\begin{equation}\label{eq:small-target-form}
 \|su_k\|_{0,k}^2
 =\langle\dbar v_k,su_k\rangle_{0,k}
 =\langle v_k,\dbar^*_{0,k}(su_k)\rangle_{0,k}
 \longrightarrow0.
\end{equation}

To retain the original class while passing to a limit, fix a smooth
metric on $L_1$ over all of $V$, with local weight $h_*$ and use the same complete
background $\omega_c$ on $U$.  Uniform upper bounds in a finite
framed cover give the global metric comparison
\begin{equation}\label{eq:fixed-hilbert-comparison}
 h_{1,k}\le h_*+C_*,\qquad
 e^{-h_*}\le e^{C_*}e^{-h_{1,k}}.
\end{equation}
It gives one bounded inclusion from each varying weighted space into
the fixed space, in both degrees zero and one.  In particular the
$u_k$ have uniformly bounded fixed norms.

On any compact subset $K\subset U$ the multiplier $s$ is nowhere
zero, and the target weights $h_{0,k}$ have uniform upper bounds.
Thus \Cref{eq:small-target-form} implies that $u_k\to0$ in the fixed
$L^2$ norm on $K$.  This local strong convergence and the global bound
imply global weak convergence to zero.  Indeed, approximate any fixed
$L^2$ test form by one supported in a compact subset of $U$; its
pairing tends to zero by local convergence, and the norm of the
remaining tail controls its pairing uniformly in $k$.  Concentration
of norm near $D$ does not affect this weak convergence.

For each $k$, choose actual exact forms $\dbar z_{k,l}$ converging to
$\gamma-u_k$ in \Cref{eq:reduced-harmonic-representative}.  By
\Cref{eq:fixed-hilbert-comparison}, their primitives belong to the
fixed degree-zero space and their derivatives converge in the fixed
degree-one norm.  Therefore $\gamma-u_k$ belongs to the fixed closure
of actual exact forms.  Apply \Cref{prop:uniform-primitives} once
more, now to the fixed smooth metric with weight $h_*$.  Its continuous class
map $\kappa_*$ annihilates that fixed closure, so
\[
 \kappa_*(u_k)=\kappa_*(\gamma)=[\gamma]
 \quad\hbox{for every }k.
\]
The fixed closed-form space is a closed Hilbert subspace, and a
continuous linear map from it to finite-dimensional ordinary
cohomology is weakly continuous.  Since $u_k\rightharpoonup0$, the
last displayed identity forces $[\gamma]=0$.

This proves analytic Dolbeault injectivity.  Dolbeault theory and GAGA
identify it with the natural map in algebraic coherent cohomology,
proving \Cref{thm:interior-injectivity}.
